\documentclass[11pt]{article}

\usepackage[utf8]{inputenc}
\usepackage[T1]{fontenc}
\usepackage{geometry}
\usepackage{amsmath}
\usepackage{amssymb}
\usepackage{booktabs}
\usepackage{hyperref}
\usepackage{xurl}
\usepackage{xcolor}
\usepackage{listings}
\usepackage{enumitem}
\usepackage{parskip}
\usepackage{graphicx}
\graphicspath{{../images/}}

\geometry{margin=1in}

\lstset{
  language=Java,
  basicstyle=\ttfamily\small,
  keywordstyle=\color{blue},
  commentstyle=\color{gray},
  stringstyle=\color{teal},
  showstringspaces=false,
  columns=fullflexible,
  frame=single,
  framerule=0.2pt,
  breaklines=true
}

\setlist[itemize]{topsep=0.3em,itemsep=0.2em}
\setlist[enumerate]{topsep=0.3em,itemsep=0.2em}

% Simple per-section source line.
\newcommand{\notesources}[1]{\par\noindent{\small\color{gray}\textit{Sources:} \sloppy #1\par}}

% Unit-wise source strings for this subject (used by slides/notes).
% Path is relative to the lecture `latex/` folder (the compilation CWD).
% OOPS with Java (CSEG1044) — unit-wise sources
%
% These are short, student-facing reference strings intended to be used with:
%   - \slidesources{...}  (in beamer slides)
%   - \notesources{...}   (in lecture notes)
%
% Style inspiration: other subjects in My-Subjects/ use semicolon-separated sources
% and include an "accessed" date for web documentation.

\newcommand{\OOPSUnitISources}{Schildt, \textit{Java: A Beginner's Guide} (10e), McGraw-Hill (Java fundamentals + OOP basics).; Oracle Java Tutorials: Getting Started + Language Basics + Classes and Objects (accessed \today).; Java SE API docs (e.g., \texttt{String}, \texttt{Scanner}) (accessed \today).}

\newcommand{\OOPSUnitIISources}{Schildt, \textit{Java: The Complete Reference} (12e), McGraw-Hill (classes, inheritance, packages, interfaces).; Bloch, \textit{Effective Java} (3e), Addison-Wesley, 2018 (design choices; interfaces vs abstract classes).; Oracle Java Tutorials: Classes and Objects + Interfaces + Packages (accessed \today).; Java SE API docs (e.g., \texttt{Comparable}, \texttt{Runnable}) (accessed \today).}

\newcommand{\OOPSUnitIIISources}{Schildt, \textit{Java: The Complete Reference} (12e) (nested/inner classes, exceptions, threads, I/O).; Oracle Java Tutorials: Nested Classes + Exceptions + Concurrency + Basic I/O (accessed \today).; Java SE API docs (e.g., \texttt{Thread}, \texttt{AutoCloseable}) (accessed \today).}

\newcommand{\OOPSUnitIVSources}{Schildt, \textit{Java: The Complete Reference} (12e) (generics, lambdas, Swing, JDBC).; Bloch, \textit{Effective Java} (3e) (generics + lambdas).; Oracle Java Tutorials: Generics + Lambda Expressions + Swing + JDBC Basics (accessed \today).; Java SE API docs (e.g., \texttt{java.util.function}, \texttt{javax.swing}, \texttt{java.sql}) (accessed \today).}

\newcommand{\OOPSUnitVSources}{Schildt, \textit{Java: The Complete Reference} (12e) (Collections Framework + wrapper classes).; Oracle Java docs: Collections Framework overview (accessed \today).; Java SE API docs (\texttt{java.util} collections; wrapper classes in \texttt{java.lang}) (accessed \today).}

\newcommand{\OOPSUnitVISources}{Git SCM: \textit{Pro Git} (2e) (version control workflow), accessed \today.; GitHub Docs (repositories, issues, pull requests), accessed \today.; Oracle Java Tutorials: Swing + JDBC (accessed \today).; JUnit 5 User Guide (accessed \today).; Apache Maven Project: Getting Started (accessed \today).; Gradle User Manual: Getting Started (accessed \today).; UML class diagrams: UML 2.x overview/spec (accessed \today).}

\newcommand{\OOPSMiscWhyInterfacesSources}{Schildt, \textit{Java: The Complete Reference} (12e) (interfaces).; Bloch, \textit{Effective Java} (3e) (prefer interfaces where appropriate).; Oracle Java Tutorials: Interfaces (accessed \today).; Java SE API docs (e.g., \texttt{Comparable}, \texttt{Runnable}) (accessed \today).}



\title{Object Oriented Programming (CSEG1044)\\Unit 01 -- Lecture 04 Notes\\Control Flow + Arrays + Strings}
\begin{document}
\maketitle
\tableofcontents

\section{Lecture Snapshot}
\begin{itemize}
  \item \textbf{Unit focus:} Introduction to OOPs and Java Fundamentals
  \item \textbf{Main new ideas:} Program control statements (selection + loops), Arrays (1D/2D), types of arrays, common patterns, Strings and \texttt{StringBuilder}
  \item \textbf{Course thread:} Use Control Flow + Arrays + Strings to make the first text-based version of Campus Resource Manager easier to reason about before the full object model arrives.
\end{itemize}
\notesources{\OOPSUnitISources}

\section{Lecture Overview}
This lecture teaches you how to \textbf{control program flow} (decisions + loops) and work with two core data structures:
\textbf{arrays} (for many values of the same type) and \textbf{strings} (text).

These three topics appear in almost every lab/program you will write.
\notesources{\OOPSUnitISources}

\section{Prerequisite Bridge}
Before reading the new material in depth, do a fast recall of the older ideas listed below.
\begin{itemize}
  \item \textbf{Types and operators}: Conditions, loop bounds, indices, and string logic depend on the previous lecture. Main reference: \texttt{Unit 01 / Lecture 03 slides/notes}.
  \item \textbf{Program flow}: Control statements only make sense once you already know the default top-to-bottom execution order. Main reference: \texttt{Unit 01 / Lecture 02 slides/notes}.
\end{itemize}
This lecture only revisits those ideas briefly so that the main explanation can stay focused on the new concept sequence.
\notesources{\OOPSUnitISources}

\section{Learning Outcomes}
After this lecture, you should be able to:
\begin{itemize}
  \item Use \texttt{if/else}, \texttt{switch}, and loops to control program flow.
  \item Create/traverse 1D and 2D arrays and solve simple array problems.
  \item Explain \texttt{String} immutability and use \texttt{StringBuilder} when needed.
\end{itemize}

\section{Suggested Reading Order}
\begin{enumerate}
  \item Read the \textbf{Prerequisite Bridge} first and confirm each recalled idea.
  \item Study the central explanation in this order: \textit{Program control statements (selection + loops)} $\rightarrow$ \textit{Arrays (1D/2D), types of arrays, common patterns} $\rightarrow$ \textit{Strings and \texttt{StringBuilder}}.
  \item Trace the worked example line by line before checking short answers.
  \item Attempt the practice section without looking at the solution immediately.
  \item Finish with the exit questions and further-study suggestions to confirm retention.
\end{enumerate}
\notesources{\OOPSUnitISources}

\section{Key Terms (Quick)}
\begin{itemize}
  \item \textbf{Selection}: choosing one path among many (\texttt{if}, \texttt{switch}).
  \item \textbf{Iteration}: repeating work until a condition changes (\texttt{for}, \texttt{while}).
  \item \textbf{Array}: fixed-size contiguous collection of same-type elements.
  \item \textbf{Index}: position in array (0-based in Java).
  \item \textbf{Immutable}: cannot be changed after creation (e.g., \texttt{String}).
  \item \textbf{StringBuilder}: mutable sequence of characters (efficient for repeated concatenation).
\end{itemize}

\section{Detailed Notes}
\subsection{Program control statements (selection + loops)}
\subsubsection{\texttt{if/else}}
Use \texttt{if} when you have conditions based on boolean expressions.

\textbf{Good practice:}
\begin{itemize}
  \item keep conditions readable (break complex conditions into variables),
  \item avoid deeply nested \texttt{if} blocks (use early returns in methods later).
\end{itemize}

\subsubsection{\texttt{switch}}
Use \texttt{switch} when one variable is compared against many constant values.
It often makes menu-driven programs readable.

\textbf{Pitfall:} forgetting \texttt{break} can cause fall-through (sometimes intended, often a bug).

\subsubsection{Loops}
\begin{itemize}
  \item \texttt{for}: when you know the number of iterations (or array traversal).
  \item \texttt{while}: when you loop until a condition changes.
  \item \texttt{do-while}: when the body must run at least once.
\end{itemize}
\textbf{Pitfall:} off-by-one errors. Remember arrays are 0-based and last index is \texttt{length - 1}.

\subsection{Arrays (1D/2D), types of arrays, common patterns}
\subsubsection{1D arrays}
Create:
\begin{itemize}
  \item \texttt{int[] a = new int[5];} (all zeros initially)
  \item \texttt{int[] b = \{1,2,3\};} (initialized)
\end{itemize}
Access:
\begin{itemize}
  \item \texttt{a[i]} for element, \texttt{a.length} for size (not a method).
\end{itemize}
Common patterns:
\begin{itemize}
  \item sum/average, max/min, linear search, counting, reversing.
\end{itemize}

\subsubsection{2D and jagged arrays}
2D array:
\begin{itemize}
  \item \texttt{int[][] m = new int[rows][cols];}
\end{itemize}
Jagged array (rows can have different lengths):
\begin{itemize}
  \item \texttt{int[][] j = new int[3][];} then allocate each row separately.
\end{itemize}
Pitfall: always check \texttt{m[i].length} for column length in jagged arrays.

\subsection{Strings and \texttt{StringBuilder}}
\subsubsection{String immutability}
In Java, \texttt{String} is immutable: operations like \texttt{toLowerCase()} or \texttt{trim()} return a new string.

\textbf{Why it matters:}
\begin{itemize}
  \item Safer sharing: if two variables reference the same string, it can't be modified unexpectedly.
  \item Performance implication: repeated concatenation in loops creates many temporary objects.
\end{itemize}

\subsubsection{\texttt{==} vs \texttt{equals}}
Use:
\begin{itemize}
  \item \texttt{==} to compare references (same object or not),
  \item \texttt{equals()} to compare string content.
\end{itemize}

\subsubsection{When to use \texttt{StringBuilder}}
If you build a string repeatedly (especially in a loop), use \texttt{StringBuilder} and \texttt{append()}.
\notesources{\OOPSUnitISources}

\section{Running Project Connection}
Use Control Flow + Arrays + Strings to make the first text-based version of Campus Resource Manager easier to reason about before the full object model arrives.
\begin{itemize}
  \item Use Campus Resource Manager as the running case study, or map the same idea to your own project domain if you are using another example.
  \item Ask where today's idea should live: model, service, DAO, UI, or team workflow.
  \item Use the lecture example as a smaller version of the same design choice you will make in the capstone.
\end{itemize}
\notesources{\OOPSUnitISources}

\section{Common Mistakes and Confusions}
\begin{itemize}
  \item Writing off-by-one loop boundaries.
  \item Comparing String values with == instead of an equality method.
  \item Treating String as mutable and expecting in-place updates.
\end{itemize}
\notesources{\OOPSUnitISources}

\section{Worked Example: Grades (control flow + arrays + strings)}
\begin{lstlisting}
import java.util.Scanner;

public class GradeStats {
    public static void main(String[] args) {
        Scanner sc = new Scanner(System.in);

        System.out.print("Enter number of students: ");
        int n = sc.nextInt();
        int[] marks = new int[n];

        for (int i = 0; i < n; i++) {
            System.out.print("Marks for student " + (i + 1) + ": ");
            marks[i] = sc.nextInt();
        }

        int max = marks[0];
        int sum = 0;
        for (int x : marks) {
            sum += x;
            if (x > max) max = x;
        }

        double avg = sum / (double) n;
        System.out.printf("Max=%d, Avg=%.2f%n", max, avg);

        String result = (avg >= 40) ? "PASS" : "NEEDS IMPROVEMENT";
        System.out.println("Class status: " + result);

        sc.close();
    }
}
\end{lstlisting}

\textbf{What to notice:}
\begin{itemize}
  \item Off-by-one handling in prompts: \texttt{i+1} for human-friendly numbering.
  \item Enhanced for-loop (\texttt{for (int x : marks)}) to traverse arrays.
  \item Casting for average: \texttt{sum / (double)n} to avoid integer division.
\end{itemize}

\section{Demo / Observation Task}
Use this block as a short reproducible demo or observation task.
\begin{itemize}
  \item Reproduce the lecture's main example around \textit{Program control statements (selection + loops)} once without looking at the solution first.
  \item Change one input, rule, or design choice in Campus Resource Manager and predict the result before checking it.
  \item Record one success case, one edge case, and one failure case so the idea becomes observable instead of purely theoretical.
\end{itemize}
\notesources{\OOPSUnitISources}

\section{Practice Check (with brief solutions)}
\begin{enumerate}
  \item What is the last valid index of an array of length 10?\; \textbf{Answer:} 9.
  \item Should we use \texttt{==} or \texttt{equals} to compare strings by content?\; \textbf{Answer:} \texttt{equals}.
  \item Which loop runs at least once: \texttt{while} or \texttt{do-while}?\; \textbf{Answer:} \texttt{do-while}.
\end{enumerate}

\section{Self-Study Practice (no solutions)}
\begin{enumerate}
  \item Write a program to find the second largest element in an array.
  \item Write a program to check if a string is a palindrome (ignore spaces and case).
  \item Create a 2D matrix and compute row-wise sums.
\end{enumerate}

\section{Further Study}
\begin{itemize}
  \item Revisit the prerequisite references if any recall bullet still feels weak.
  \item Rework the lecture example with one changed assumption, input, or design choice.
  \item Read the textbook and documentation references listed in the source line for a second explanation of the same topic.
  \item Apply the same idea once inside Campus Resource Manager so that the concept moves from theory to design practice.
\end{itemize}
\notesources{\OOPSUnitISources}

\section{Exit Questions (with sample answers)}
\textbf{Q1.} Why is \texttt{String} immutable and why does it matter?\par
\textbf{Sample answer:} Immutability makes strings safe to share and helps with security and caching. It also means operations create new objects, so repeated concatenation in loops can be inefficient.

\vspace{0.6em}
\textbf{Q2.} How do you avoid off-by-one errors in array loops?\par
\textbf{Sample answer:} Remember Java arrays are 0-based and the last index is \texttt{length - 1}. Use loop conditions like \texttt{i < a.length} rather than \texttt{i <= a.length}.

\vspace{0.6em}
\textbf{Q3.} Why is String immutability useful?\par
\textbf{Sample answer:} It makes String values safer to share and easier to reason about across APIs and larger programs.

\end{document}
